\section{Tok manbalarini ketma-ket va parallel ulash}



\begin{enumerate}

\item Har birining EYKi $6$ V va ichki qarshiligi $0,6~\Omega$ bo‘lgan ikkita akkumulyator ketma-ket ulanganda, $2$ A tok bo‘lishi uchun tashqi zanjirdagi qarshilik qanday bo‘lishi kerak ($\Omega$)?
\begin{enumerate}[label=\Alph*)]
\item 48
\item 18
\item 4,8
\item 2
\end{enumerate}

\item Tashqi zanjirdagi kuchlanish $100$ V va tok kuchi $5$ A bo‘lishi uchun EYKi $2,5$ V dan va ichki qarshiligi $0,1~\Omega$ dan bo‘lgan manbalardan nechtasini ketma-ket ulab, batareya tuzish kerak bo‘ladi?
\begin{enumerate}[label=\Alph*)]
\item 25
\item 36
\item 40
\item 50
\end{enumerate}

\item Rasmda ko‘rsatilgan elektr zanjirda $\varepsilon_1=6$ V, $\varepsilon_2=18$ V, $R=3~\Omega$. Manbalar ichki qarshilikka ega emas. Zanjirdagi tok kuchi necha Amper?

\begin{circuitikz}[european]
\draw
  % Chap tomondagi vertikal rezistor R
  (0,0) to[R=R] (0,3)
  
  % Yuqori tomondagi tok manbai epsilon_1
  % Musbat qutbi o'ng tomonda
  (0,3) to[battery1, l=$\varepsilon_1$] (4,3)
  
  % O'ng tomondagi vertikal rezistor R
  (4,3) to[R=R] (4,0)
  
  % Pastki tomondagi tok manbai epsilon_2
  % Musbat qutbi chap tomonda
  (4,0) to[battery1, l=$\varepsilon_2$] (0,0);
\end{circuitikz}

\begin{enumerate}[label=\Alph*)]
\item 8
\item 6
\item 4
\item 3
\end{enumerate}

\item $\varepsilon_1=5$ V, $\varepsilon_2=3$ V, $\varepsilon_3=4$ V, $r=1~\Omega$ va $R=7~\Omega$ bo‘lsa, $R$ qarshilikdagi kuchlanishning tushishi necha Volt bo‘ladi (rasmga q.)?

\begin{circuitikz}[european]
\draw
  % Pastki tarmoq: Tashqi rezistor R
  (0,0) to[R=R] (4,0)
  
  % O'ng tarmoq: Tok manbai epsilon_3 va ichki qarshilik r
  % Qutblari: yuqorida +, pastda -
  (4,0) to[battery1, l_=$\varepsilon_3$, i^=$r$] (4,3)
  
  % Yuqori tarmoq: Tok manbai epsilon_2 va ichki qarshilik r
  % Qutblari: chapda +, o'ngda -
  (4,3) to[battery1, l=$\varepsilon_2$, i=$r$] (0,3)
  
  % Chap tarmoq: Tok manbai epsilon_1 va ichki qarshilik r
  % Qutblari: yuqorida +, pastda -
  (0,0) to[battery1, l=$\varepsilon_1$, i=$r$] (0,3);
\end{circuitikz}

\begin{enumerate}[label=\Alph*)]
\item 14
\item 1,4
\item 7
\item 0,7
\end{enumerate}

\item Rasmda ko‘rsatilgan elektr zanjirda $\varepsilon_1=6$ V, $\varepsilon_2=12$ V va $R=6~\Omega$. Manbalarning ichki qarshiligi bo‘lmasa, zanjirdagi tok kuchi necha Amper bo‘ladi?

\begin{circuitikz}[european]
\draw
  % Chap vertical tarmoq: epsilon_1 manbasi
  % Musbat qutbi (+) pastda joylashgan
  (0,0) to[battery1, l=$\varepsilon_1$] (0,3)
  
  % Yuqori horizontal tarmoq: R rezistor
  (0,3) -- (1.5,3) to[R=R] (2.5,3) -- (4,3)
  
  % O'ng vertical tarmoq: epsilon_2 manbasi
  % Musbat qutbi (+) pastda joylashgan
  (4,3) to[battery1, l_=$\varepsilon_2$] (4,0)
  
  % Pastki horizontal tutashtiruvchi sim
  (4,0) -- (0,0);
\end{circuitikz}

\begin{enumerate}[label=\Alph*)]
\item 1
\item 2
\item 3
\item 4
\end{enumerate}

\item Rasmda ko‘rsatilgan zanjirdagi tok kuchi necha Amperga teng? $\varepsilon_1=\varepsilon_2=\varepsilon_3=6$ V, $r_1=r_2=r_3=0,5~\Omega, R=1,5~\Omega$.

\begin{circuitikz}[european]
\draw
  % Yuqori tarmoq: epsilon_1 va epsilon_2 manbalari
  % epsilon_1 musbat qutbi o'ngda, epsilon_2 musbat qutbi chapda
  (0,3) to[battery1, l=$\varepsilon_1$] (2,3)
        to[battery1, l=$\varepsilon_2$, invert] (4,3)
  
  % O'ng tarmoq: R rezistor
  (4,3) to[R=R] (4,0)
  
  % Pastki tarmoq: epsilon_3 manbasi
  % Musbat qutbi chapda joylashgan
  (4,0) to[battery1, l_=$\varepsilon_3$] (0,0)
  
  % Chap tarmoq: tutashtiruvchi sim
  (0,0) -- (0,3);
\end{circuitikz}

\begin{enumerate}[label=\Alph*)]
\item 1,5
\item 2
\item 3
\item 6
\end{enumerate}

\item Uchta tok manbai va bitta qarshilikdan rasmda ko‘rsatilgandek zanjir tuzilgan. Agar $\varepsilon_1=3,5$ V, $\varepsilon_2=1,5$ V, $\varepsilon_3=2$ V, $r_1=r_2=r_3=0,2~\Omega$, va $R=4,4~\Omega$ bo‘lsa, $R$ qarshilikdan o‘tayotgan tok kuchi necha Amper?

\begin{circuitikz}[european]
\draw
  % Yuqori tarmoq: epsilon_1 va epsilon_2 manbalari
  % epsilon_1 (+) o'ngda, epsilon_2 (+) chapda (manbadagi rasmga ko'ra)
  (0,3) to[battery1, l=$\varepsilon_1$] (2,3)
        to[battery1, l=$\varepsilon_2$, invert] (4,3)
  
  % O'ng vertical tarmoq: Tashqi rezistor R
  (4,3) to[R=R] (4,0)
  
  % Pastki tarmoq: epsilon_3 manbasi
  % (+) chapda joylashgan (4,0 dan 0,0 ga qarab)
  (4,0) to[battery1, l_=$\varepsilon_3$] (0,0)
  
  % Chap vertical tarmoq: tutashtiruvchi sim
  (0,0) -- (0,3);
\end{circuitikz}

\begin{enumerate}[label=\Alph*)]
\item 0,6
\item 0,9
\item 1,6
\item 0,8
\end{enumerate}

\item Sxemada $\varepsilon_1=2$ V, $\varepsilon_2=1$ V, $r_1=r_2=1~\Omega, R=2~\Omega$. Tashqi qarshilikdagi kuchlanishning tushishini toping (V).

\begin{circuitikz}[european]
\draw
  % Yuqori horizontal tarmoq: Ikkita tok manbai ketma-ket
  % Har ikkala manbaning musbat qutbi (+) o'ng tomonga yo'nalgan
  (0,3) to[battery1, l=$\varepsilon_1$, annotation=$r_1$] (2,3)
        to[battery1, l=$\varepsilon_2$, annotation=$r_2$] (4,3)
  
  % O'ng vertical tutashtiruvchi sim
  (4,3) -- (4,0)
  
  % Pastki horizontal tarmoq: Tashqi rezistor R
  (4,0) to[R=R] (0,0)
  
  % Chap vertical tutashtiruvchi sim
  (0,0) -- (0,3);
\end{circuitikz}

\begin{enumerate}[label=\Alph*)]
\item 0,5
\item 1
\item 1,5
\item 2
\end{enumerate}

\item Kondensatorlardagi kuchlanishlar qanday (V)? $\varepsilon_1=12$ V; $\varepsilon_2=13$ V; $C_1=3~\mu$F; $C_2=7~\mu$F. 

\begin{circuitikz}[european]
\draw
  % Yuqori horizontal tarmoq: epsilon_1 manbasi
  % Musbat qutbi o'ng tomonda
  (0,3) to[battery1, l=$\varepsilon_1$] (4,3)
  
  % O'ng vertical tarmoq: Kondensator C2
  (4,3) to[C, l=$C_2$] (4,0)
  
  % Pastki horizontal tarmoq: epsilon_2 manbasi
  % Musbat qutbi chap tomonda
  (4,0) to[battery1, l_=$\varepsilon_2$] (0,0)
  
  % Chap vertical tarmoq: Kondensator C1
  (0,0) to[C, l=$C_1$] (0,3);
\end{circuitikz}

\begin{enumerate}[label=\Alph*)]
\item 6,5; 18,5
\item 7; 18
\item 18; 7
\item 17,5; 7,5
\end{enumerate}

\item Sxemada EYK lar $\varepsilon_1=3$ V va $\varepsilon_2=2$ V, ichki qarshiliklar $r=0,5~\Omega$ bo‘lsa, voltmetr necha Voltni ko‘rsatadi? $R_V = \infty$.

\begin{circuitikz}[european]
\draw
  % Yuqori tarmoq: epsilon_1 va r
  % Musbat qutb o'ng tomonda
  (0,2) to[battery1, l=$\varepsilon_1$] (1.5,2) to[R=r] (3,2)
  
  % O'rta tarmoq: epsilon_2 va r
  % Musbat qutb o'ng tomonda
  (0,1) to[battery1, l=$\varepsilon_2$] (1.5,1) to[R=r] (3,1)
  
  % Chap va o'ng vertikal tutashtiruvchi simlar
  (0,2) -- (0,1) -- (0,0)
  (3,2) -- (3,1) -- (3,0)
  
  % Pastki qism: Voltmetr (V)
  (0,0) to[voltmeter, l=V] (3,0);
\end{circuitikz}

\begin{enumerate}[label=\Alph*)]
\item 1
\item 2,5
\item 2
\item 3
\end{enumerate}

\item Sxemada EYK lar $\varepsilon_1=3$ V va $\varepsilon_2=2$ V, ichki qarshiliklar $r=0,5~\Omega$ bo‘lsa, voltmetr necha Voltni ko‘rsatadi? $R_V = \infty$. 

\begin{circuitikz}[european]
\draw
  % Yuqori tarmoq: epsilon_1 va r
  % Musbat qutb o'ng tomonda
  (0,2) to[battery1, l=$\varepsilon_1$] (1.5,2) to[R=r] (3,2)
  
  % O'rta tarmoq: epsilon_2 va r
  % Musbat qutb chap tomonda (invert orqali)
  (0,1) to[battery1, l=$\varepsilon_2$, invert] (1.5,1) to[R=r] (3,1)
  
  % Chap va o'ng vertikal tutashtiruvchi simlar
  (0,2) -- (0,1) -- (0,0)
  (3,2) -- (3,1) -- (3,0)
  
  % Pastki qism: Voltmetr (V)
  (0,0) to[voltmeter, l=V] (3,0);
\end{circuitikz}

\begin{enumerate}[label=\Alph*)]
\item 0,5
\item 1
\item 2
\item 3
\end{enumerate}

\item Elektr yurituvchi kuchlari teng, ichki qarshiliklari $2$ marta farq qiladigan tok manbalari rasmdagidek ulangan. Voltmetrning ko‘rsatishini toping. $R_V = \infty$.

\begin{circuitikz}[european]
\draw
  % Chap tarmoq: epsilon va r
  (0,0) to[battery1, l=$\varepsilon$] (0,1.5)
        to[R=r] (0,3)
  
  % O'rta tarmoq: Voltmetr (V)
  (2,0) to[voltmeter, l=V, *-*] (2,3)
  
  % O'ng tarmoq: epsilon va 2r
  (4,0) to[battery1, l_=$\varepsilon$] (4,1.5)
        to[R=2r] (4,3)
  
  % Yuqori va pastki horizontal tutashtiruvchi simlar
  (0,3) -- (4,3)
  (0,0) -- (4,0);
\end{circuitikz}

\begin{enumerate}[label=\Alph*)]
\item 0
\item $\varepsilon/2$
\item $\varepsilon$
\item $\varepsilon/3$
\end{enumerate}

\item EYUK $1,5$ V va $2$ V bo‘lgan elementlar bir xil ishorali qutblari bilan ulangan. Batareya klemmalariga ulangan voltmetr $1,7$ V kuchlanishni ko‘rsatdi. Elementlar ichki qarshiliklari nisbatini toping.
\begin{enumerate}[label=\Alph*)]
\item 1
\item 2/3
\item 2
\item 1/3
\end{enumerate}

\item EYK lari $\varepsilon_1$ va $\varepsilon_2$ bo‘lgan manbalar bir xil qutblari bilan o‘zaro parallel ulangan. Agar manbalarning ichki qarshiliklari teng bo‘lsa, manba qutblaridagi potensiallar farqi qanday?
\begin{enumerate}[label=\Alph*)]
\item $\varepsilon_1 + \varepsilon_2$
\item $(\varepsilon_1 - \varepsilon_2)/2$
\item 0
\item $(\varepsilon_1 + \varepsilon_2)/2$
\end{enumerate}

\end{enumerate}



\begin{table}[!ht]
\centering
\begin{tabular}{|c|c|c|c|c|c|c|c|c|c|}
\hline
1d & 2d & 3d & 4b & 5c & 6d & 7c & 8b & 9c & 10a \\ \hline
11b & 12d & 13c & 14d &  &  &  &  &  &  \\ \hline
\end{tabular}
\end{table}


